\chapter{Introduction}
\label{chap:introduction}

% ─────────────────────────────────────────────────────────────────────────────
\section{Motivation}
\label{sec:motivation}
% ─────────────────────────────────────────────────────────────────────────────

This section situates the thesis within the broader challenge of clinical decision-making
and explains why multi-agent systems represent the most promising architectural
direction for clinical AI.  It traces the progression from single-model limitations through
the emergence of coordinated multi-agent designs, and identifies the specific unsolved
challenge that this thesis addresses.

Clinical decision-making is one of the most cognitively demanding tasks in modern
medicine.  At every stage of patient care --- from initial triage and history-taking
through investigation ordering, differential diagnosis, and treatment planning ---
clinicians must synthesise information from multiple heterogeneous sources
simultaneously, weigh competing hypotheses under conditions of uncertainty, and act
under time pressure.  The consequences of errors at any stage can be severe.  Despite
decades of investment in clinical decision support systems, the reality of most hospital
environments is that physicians still navigate these tasks largely in isolation, relying
on individual memory, pattern recognition, and peer consultation when circumstances
allow~\cite{Fan2024AIHospital}.

Artificial Intelligence (AI) has long been proposed as a tool to augment clinical
decision-making, and the emergence of large language models (LLMs) has significantly
renewed this ambition.  Unlike earlier rule-based clinical decision support tools, LLMs
trained on vast corpora of medical text are capable of generating contextually relevant
clinical reasoning, synthesising information across specialties, and communicating
findings in natural language~\cite{Abbasian2024openCHA}.  They have demonstrated
impressive performance on standardised medical examinations and benchmark
question-answering tasks~\cite{Tang2024MedAgents}.  Yet, despite this apparent
capability, their deployment in real clinical workflows has lagged substantially behind
the expectations raised by laboratory evaluations~\cite{Yan2024ClinicalLab}.

The fundamental limitation is architectural: a single LLM, regardless of its scale or
training quality, operates as a monolithic system that processes all clinical inputs
through a single forward pass, produces a single output, and has no capacity to verify
its own reasoning, consult an external knowledge source, or correct an error it has
already made.  Tang~et~al.~\cite{Tang2024MedAgents} demonstrated empirically that 77\%
of errors made by LLMs on medical tasks are attributable not to failures of reasoning
but to domain knowledge gaps.  Ke~et~al.~\cite{Ke2024CognitiveBias} showed that single
LLMs systematically exhibit the same cognitive biases --- anchoring, confirmation, and
availability bias --- that are the primary sources of diagnostic error in human
clinicians, and that these biases cannot be self-corrected within a single-agent
architecture.  Shang~et~al.~\cite{Shang2025DynamiCare} further demonstrated that
single-pass reasoning is fundamentally inadequate for complex clinical cases, where
dynamic multi-round information gathering significantly outperforms static approaches.
Bedi~et~al.~\cite{Bedi2025OptParadox} demonstrated on 2{,}400 real patient cases that
even the best-performing individual LLM agents, when assembled into a system without
careful coordination design, produce a hallucination rate of 13.87\% on test result
reporting --- a level that would be considered clinically unacceptable in any deployed
system.  Choi~et~al.~\cite{Choi2024MALADE} corroborated this concern in the
pharmacovigilance domain, showing that non-RAG single-agent systems produce
substantially higher rates of hallucinated drug--event associations than their
multi-agent counterparts.

Multi-agent systems (MAS) offer a structural solution to these limitations.  Rather than
placing the full burden of clinical reasoning on a single model, a multi-agent system
distributes the task across a coordinated ensemble of specialised agents --- each with a
defined role, a distinct knowledge focus, and a structured communication interface with
the other agents in the system~\cite{Kim2024MDAgents}.  The collective output is not the
output of any single agent but the product of a coordinated workflow in which agents
check each other's work, challenge each other's conclusions, and escalate unresolved
disagreements to a higher authority.  This architectural shift from monolithic to
multi-agent design mirrors the shift from individual clinical decision-making to
multidisciplinary team (MDT) consultation that clinical medicine has already adopted as
best practice for complex
cases~\cite{Fan2024AIHospital,MAC2025npj}.  Chen~et~al.~\cite{Chen2025MDTeamGPT} showed
that simulating such MDT consultations through structured multi-agent debate produces
higher diagnostic accuracy than any single-agent configuration, particularly when agents
begin deliberation with heterogeneous initial diagnoses.

The case for multi-agent systems in clinical AI is no longer merely theoretical.  Recent
works --- MDAgents at NeurIPS 2024~\cite{Kim2024MDAgents}, ZODIAC at cardiologist-level
performance~\cite{Zhou2024ZODIAC}, ClinicalLab across eleven real hospital
departments~\cite{Yan2024ClinicalLab}, and the MAC Framework published in \textit{npj
Digital Medicine}~\cite{MAC2025npj} --- have demonstrated that multi-agent LLM
frameworks can approach or match clinician-level performance in specific clinical tasks.
Chen~et~al.~\cite{Chen2024RareAgents} demonstrated that even an open-source model
(Llama-3.1-70B) can outperform GPT-4o on rare disease diagnosis when embedded in a
well-designed multi-agent architecture with persistent memory and tool-use capabilities,
underscoring that coordination architecture matters more than individual model scale.
However, a critical unresolved challenge remains: the design of multi-agent clinical
workflows that are not only accurate but also robust, safe, and generalisable across the
full complexity of real clinical scenarios.  Bedi~et~al.~\cite{Bedi2025OptParadox}
identified what they term the Optimisation Paradox --- the finding that a system
assembled from individually best-performing components consistently underperforms a less
capable but better-coordinated system --- establishing workflow coordination architecture
as the primary unsolved design challenge in clinical AI.  It is this challenge that this
thesis addresses.

% ─────────────────────────────────────────────────────────────────────────────
\section{Problem Statement}
\label{sec:problem}
% ─────────────────────────────────────────────────────────────────────────────

Having established the motivation for multi-agent clinical AI, this section identifies the
four specific engineering challenges that existing frameworks fail to address
simultaneously, and frames the precise research gap that this thesis
targets.

Despite the substantial progress documented in the recent multi-agent clinical AI
literature, the field currently lacks a validated, end-to-end multi-agent system for
clinical decisioning that simultaneously addresses the four inter-related engineering
challenges that real-world clinical deployment requires: coordinated workflow
orchestration across heterogeneous data sources, adaptive conflict resolution between
disagreeing agents, automated error recovery within the pipeline, and evaluable safety
properties~\cite{Kim2024MDAgents,Kim2025TAO,Wu2024EHRFlow}.

The \textbf{first challenge is coordination}.  As Bedi~et~al.~\cite{Bedi2025OptParadox}
demonstrated, assembling individually capable agents into a system without careful
coordination design produces worse outcomes than a less capable but coherently designed
system.  Existing frameworks address coordination in narrow ways: MDAgents~\cite{Kim2024MDAgents}
uses complexity-based routing but cannot replan once execution has begun; TAO~\cite{Kim2025TAO}
uses tiered escalation for safety but was not designed for end-to-end diagnostic
workflows; EHRFlow~\cite{Wu2024EHRFlow} implements sequential error recovery but
operates only on structured EHR analysis rather than open-ended clinical decisioning.
No existing framework provides a generalisable coordination architecture for the full
clinical decisioning pipeline --- from patient data intake to differential diagnosis and
management recommendation.

The \textbf{second challenge is conflict resolution}.  Fan~et~al.~\cite{Fan2024AIHospital}
demonstrated empirically that removing structured conflict resolution from a multi-agent
clinical system causes a statistically significant reduction in diagnostic accuracy,
establishing it as a structural necessity rather than an optional enhancement.  Yet
Wang~et~al.~\cite{Wang2025Catfish} identified silent agreement --- the failure of agents
to genuinely challenge one another even when evidence is ambiguous --- as a systematic
failure mode that undermines the quality of agent deliberation in most existing
frameworks.  The Catfish Agent~\cite{Wang2025Catfish} and Layered Debate
MAS~\cite{Zhao2025LayeredDebate} each address this problem in narrow experimental
settings, but neither provides a conflict resolution mechanism that is adaptive to case
complexity, integrated into a full clinical workflow, and generalisable across clinical
domains.

The \textbf{third challenge is safety}.  Chen~et~al.~\cite{Chen2025MedSentry}
demonstrated that MAS topology is a primary determinant of clinical safety and that a
single compromised agent in a centralised system can corrupt system-wide outputs ---
including generating false prescriptions and distorting diagnoses --- in ways that bypass
naive safety filters.  Only three of the twenty-two papers reviewed in Chapter~\ref{chap:background}
address safety as a primary research question, and none evaluates the quantitative
trade-off between safety constraint enforcement and diagnostic
accuracy~\cite{Kim2025TAO,Wu2024EHRFlow,Chen2025MedSentry}.

The \textbf{fourth challenge is evaluation}.  The majority of existing frameworks are
evaluated exclusively on structured multiple-choice question-answering
benchmarks~\cite{Tang2024MedAgents,Kim2024MDAgents}.  These evaluations measure whether
agents can select the correct answer from a closed set of options, which is a
fundamentally different task from generating a complete clinical management plan,
ordering appropriate investigations, and justifying a differential diagnosis from
unstructured patient data.  CDR-Agent~\cite{Xiang2025CDRAgent} and
ClinicalLab~\cite{Yan2024ClinicalLab} represent exceptions, but both evaluate
domain-specific tasks.  A rigorous evaluation framework for multi-agent clinical
decisioning must measure performance on the full task, not a simplified proxy.

This thesis addresses all four challenges through the design, implementation, and
evaluation of a multi-agent system for AI clinical decisioning via automation workflows,
built on a clinical data lake of 10{,}000 synthetic patients, transformed through a
medallion lakehouse architecture (Bronze, Silver via OMOP CDM~v5.4, Gold), and evaluated
using a comprehensive framework combining automated metrics, LLM-as-judge rationale
scoring, and ground-truth comparison~\cite{Bedi2025OptParadox,Chen2025MedSentry}.

% ─────────────────────────────────────────────────────────────────────────────
\section{Objectives}
\label{sec:objectives}
% ─────────────────────────────────────────────────────────────────────────────

This section translates the four challenges identified in the problem statement into
concrete, measurable engineering objectives that collectively define the scope of work
for this thesis.

The primary goal of this thesis is to design, implement, and evaluate an end-to-end
multi-agent system for AI clinical decisioning that operates over a structured clinical
data lake, demonstrating that coordinated multi-agent workflows outperform single-agent
baselines on clinical decisioning tasks while maintaining defined safety and reliability
properties.  The specific objectives are as follows.

\begin{itemize}

  \item \textbf{Design and implement a clinical data lake pipeline} using the Synthea
  patient generator, Python-based Bronze ingestion with FHIR~R4 parsing, dbt-core Silver
  transformation to the OMOP CDM~v5.4 standard, and a Python Gold assembler that
  produces structured EHR and laboratory case objects for agent consumption --- all
  orchestrated by a Prefect DAG with automated monitoring and retry logic.

  \item \textbf{Develop a multi-agent decisioning framework} comprising a specialised EHR
  Agent and Lab Agent built on LangChain and Claude Sonnet, coordinated by a custom
  Python Multi-Agent Orchestrator that implements adaptive complexity routing, structured
  inter-agent communication protocols, and a built-in conflict resolution mechanism for
  handling agent disagreements.

  \item \textbf{Implement adaptive workflow orchestration} in which the Orchestrator
  dynamically assigns single-agent or collaborative multi-agent processing to each
  clinical case based on assessed complexity --- mirroring the MDAgents
  and TAO architectural precedents --- and can trigger escalation and
  replanning in response to low-confidence or conflicting agent outputs.

  \item \textbf{Integrate Retrieval-Augmented Generation (RAG)} using a FAISS vector
  store built from curated clinical guideline documents, enabling agents to ground their
  reasoning in current evidence-based protocols at inference time and reducing
  hallucinated clinical recommendations --- following the demonstrated 64\% hallucination
  reduction achieved by Choi~et~al.\ in a comparable RAG-integrated
  multi-agent system.

  \item \textbf{Build and validate a data quality assurance layer} using dbt schema
  tests and Great Expectations statistical profile checks, ensuring that the clinical
  data lake on which agents operate meets defined quality standards for completeness,
  referential integrity, and clinical plausibility before any agent query is processed.

  \item \textbf{Design and execute a comprehensive evaluation framework} comparing
  single-agent and multi-agent system configurations across diagnostic accuracy,
  recommendation completeness, inter-run consistency, hallucination rate, and safety
  boundary adherence on a held-out set of 100 clinically diverse cases.

  \item \textbf{Implement an end-to-end automated error recovery mechanism} within the
  data pipeline and agent workflow, following the Debugger Agent pattern established by
  Wu~et~al.\ in EHRFlow, to ensure that execution failures do not
  propagate to downstream agents as corrupted inputs and that the system degrades
  gracefully under data quality failures.

  \item \textbf{Develop a clinical portal using Streamlit} that reads directly from the
  DuckDB clinical data lake, enabling structured review of agent outputs, data quality
  dashboard monitoring, and case-level comparison of single-agent versus multi-agent
  decisioning.

\end{itemize}

% ─────────────────────────────────────────────────────────────────────────────
\section{Thesis Outline}
\label{sec:outline}
% ─────────────────────────────────────────────────────────────────────────────

This section provides a roadmap of the thesis structure, summarising the purpose and
content of each chapter to guide the reader through the document.  The structure follows
a standard empirical research progression from literature review through system design,
experimentation, and reflection.

This thesis is structured into five chapters, each addressing a distinct phase of the
research from literature foundation through implementation and evaluation.

\textbf{Chapter~\ref{chap:background}: Background} reviews the existing literature on
multi-agent LLM frameworks for clinical AI.  Twenty-two papers are reviewed and
organised into five thematic categories: foundational frameworks, coordination
architectures, inter-agent deliberation, domain-specific deployments, and safety and
evaluation.  The chapter identifies the four research gaps that this thesis addresses and
establishes the architectural vocabulary and design precedents that inform the system
design.

\textbf{Chapter~3: Methodology} describes the complete system design, including the
clinical data lake architecture (Synthea data generation, Bronze/Silver/Gold medallion
pipeline, OMOP CDM~v5.4 transformation, dbt-core models, Prefect orchestration), the
multi-agent framework design (EHR Agent, Lab Agent, Orchestrator, RAG layer, conflict
resolution mechanism), the data quality assurance layer, and the evaluation framework
design.

\textbf{Chapter~4: Results} presents and analyses the system's performance across all
evaluation dimensions.  This includes data pipeline performance metrics, single-agent
versus multi-agent comparative results (diagnostic accuracy, hallucination rate,
consistency, recommendation completeness), conflict resolution mechanism effectiveness,
and clinical portal demonstration outputs.  Results are interpreted in the context of the
literature benchmarks established in Chapter~\ref{chap:background}.

\textbf{Chapter~5: Conclusion and Future Work} summarises the thesis contributions,
reflects on the degree to which the stated objectives were achieved, and discusses the
limitations of the study design.  Directions for future work are proposed, including
extension to real EHR data, integration of additional agent specialisms, and prospective
clinical validation with practising clinicians.
